\clearpage\thispagestyle{empty}
\pdfbookmark[0]{Quellen zu Vesta (13. August 1814)}{vesta1814aug-sources}
\begin{center}\textbf{Quellen zu Vesta (13. August 1814)}\end{center}
\noindent Diese KI-Übertragung ins heutige Deutsch folgt dem Bericht mit der Überschrift Vesta in den \href{https://gdz.sub.uni-goettingen.de/id/PPN236007467}{\textit{Werken}, Band~VI} von Gauß (Göttingen, 1874), S.~377--379. Der Bericht erschien unter Göttingen in den \textit{Göttingischen gelehrten Anzeigen}, Stück~129 (13. August 1814), S.~\href{https://images.p.sub.uni-goettingen.de/iiif/image/gdz:PPN319721507_1814:00001287/full/1600,/0/default.jpg}{1281}, \href{https://images.p.sub.uni-goettingen.de/iiif/image/gdz:PPN319721507_1814:00001288/full/1600,/0/default.jpg}{1282}, \href{https://images.p.sub.uni-goettingen.de/iiif/image/gdz:PPN319721507_1814:00001289/full/1600,/0/default.jpg}{1283}, \href{https://images.p.sub.uni-goettingen.de/iiif/image/gdz:PPN319721507_1814:00001290/full/1600,/0/default.jpg}{1284}. Die Links führen zum digitalisierten Exemplar der Niedersächsischen Staats- und Universitätsbibliothek Göttingen.
\par\smallskip\noindent Die Werkausgabe gibt für die Polhöhe der neuen Königsberger Sternwarte $54^\circ\,42'\,50''31$ an. In der vorletzten Zeile der Vergleichstabelle steht bei Nicolai neben der Rektaszensionsdifferenz $+5.58$ der 26. Februar. Der Erstdruck gibt auf S.~1282 die Polhöhe $54^\circ\,17'\,50''31$ an und druckt auf S.~1283 in dieser Zeile den Tag 29 und die Differenz $+5{,}58$. In einer \href{https://images.p.sub.uni-goettingen.de/iiif/image/gdz:PPN319721507_1814:00001996/full/1600,/0/default.jpg}{Fußnote auf S.~1988} bezeichnet das Stück vom 12. Dezember 1814 die 17 Minuten der Polhöhe als Druckfehler und gibt 42 Minuten an. KI-Übertragung und KI-Satz.
\clearpage
